\documentclass[sigconf,review,anonymous]{acmart}

%% BibTeX command
\AtBeginDocument{%
  \providecommand\BibTeX{{%
    Bib\TeX}}}

%% Packages (booktabs is already loaded by acmart)
\usepackage{subcaption}
\usepackage{listings}
\usepackage{xcolor}
\usepackage{algorithm}
\usepackage{algpseudocode}

%% Solidity language definition
\lstdefinelanguage{Solidity}{
  keywords={contract, function, modifier, public, private, external, internal, view, pure, payable, returns, return, if, else, for, while, require, assert, mapping, address, uint256, uint, bool, string, memory, storage, calldata, event, emit, struct, enum, constructor, msg, this},
  keywordstyle=\color{blue}\bfseries,
  sensitive=true,
  comment=[l]{//},
  morecomment=[s]{/*}{*/},
  commentstyle=\color{gray}\ttfamily,
  stringstyle=\color{orange}\ttfamily,
  morestring=[b]',
  morestring=[b]"
}

\lstset{
  language=Solidity,
  basicstyle=\footnotesize\ttfamily,
  backgroundcolor=\color[RGB]{248,248,248},
  numbers=left,
  numberstyle=\tiny\color{gray},
  numbersep=5pt,
  breaklines=true,
  frame=single,
  captionpos=b
}

%% Rights management
\setcopyright{acmcopyright}
\copyrightyear{2025}
\acmYear{2025}
\acmDOI{XXXXXXX.XXXXXXX}

%% Conference information (placeholder - update when accepted)
\acmConference[ASE '25]{40th IEEE/ACM International Conference on Automated Software Engineering}{September 2025}{Sacramento, CA, USA}
\acmISBN{978-1-4503-XXXX-X/25/09}

%% Submission ID (uncomment and fill when submitting)
%%\acmSubmissionID{123-A56-BU3}

%%
%% end of the preamble
%%

\begin{document}

\title{IntentKeeper: LLM-Assisted Intent Annotation and Program Repair for Solidity Smart Contracts}

\author{Anonymous Author(s)}
\affiliation{%
  \institution{Anonymous Institution}
  \city{Anonymous City}
  \country{Anonymous Country}
}
\email{anonymous@example.com}

\begin{abstract}
% TODO: Write abstract
% Key points to cover:
% - Problem: Intent annotations require manual effort; fixing violated intents is challenging
% - Solution: LLM-assisted annotation suggestion and program repair
% - Contribution: (1) Annotation suggestion, (2) Program repair based on validation results
% - Results: TBD
\end{abstract}

\begin{CCSXML}
<ccs2012>
   <concept>
       <concept_id>10011007.10011006.10011008.10011009.10011012</concept_id>
       <concept_desc>Software and its engineering~Software verification and validation</concept_desc>
       <concept_significance>500</concept_significance>
   </concept>
   <concept>
       <concept_id>10011007.10011006.10011039</concept_id>
       <concept_desc>Software and its engineering~Formal language definitions</concept_desc>
       <concept_significance>300</concept_significance>
   </concept>
</ccs2012>
\end{CCSXML}

\ccsdesc[500]{Software and its engineering~Software verification and validation}
\ccsdesc[300]{Software and its engineering~Formal language definitions}

\keywords{Smart Contracts, Solidity, Large Language Models, Intent Annotation, Program Repair, Developer Tools}

\maketitle

%% ============================================================
\section{Introduction}\label{sec:introduction}
%% ============================================================

% TODO: Motivation
% - Logic errors are hard to detect without knowing developer intent
% - IntentChecker enables intent specification but requires manual annotation
% - Two gaps: (1) annotation burden, (2) no automated repair

% TODO: Our approach
% - IntentKeeper: LLM-based extension of IntentChecker
% - Annotation Suggestion: LLM recommends appropriate intent annotations
% - Program Repair: LLM fixes code when intent is violated/uncertain

% TODO: Contributions
% - Annotation suggestion using LLM
% - Program repair guided by validation results
% - Evaluation on real-world contracts

%% ============================================================
\section{Background}\label{sec:background}
%% ============================================================

\subsection{Intent-Based Validation in Smart Contracts}
% TODO: Brief overview of IntentChecker
% - @During and @Post annotations
% - Interval-domain abstract interpretation
% - Validation results: Satisfied, Violated, Uncertain

\subsection{Large Language Models for Code}
% TODO: Overview of LLMs for code generation/repair
% - Code LLMs (CodeLlama, StarCoder, DeepSeek-Coder, etc.)
% - Fine-tuning approaches
% - Prompting strategies

%% ============================================================
\section{Approach}\label{sec:approach}
%% ============================================================

\subsection{Overview}
% TODO: System architecture of IntentKeeper
% - Integration with IntentChecker
% - LLM component

\subsection{Annotation Suggestion}\label{sec:annotation-suggestion}
% TODO: How LLM suggests intent annotations
% - Input: Solidity code (function context)
% - Output: Suggested @During/@Post annotations
% - Prompting or fine-tuning strategy

\subsection{Program Repair}\label{sec:program-repair}
% TODO: How LLM repairs code based on validation results
% - Input: Code + Intent annotation + Validation result (Violated/Uncertain) + Violated region
% - Output: Repaired code that satisfies the intent
% - Repair strategy

%% ============================================================
\section{Implementation}\label{sec:implementation}
%% ============================================================

% TODO: Implementation details
% - LLM selection and configuration
% - Fine-tuning dataset (if applicable)
% - Integration with IntentChecker

%% ============================================================
\section{Evaluation}\label{sec:evaluation}
%% ============================================================

\subsection{Research Questions}
% RQ1: How effective is annotation suggestion?
% RQ2: How effective is program repair?
% RQ3: How do different LLMs/strategies compare?

\subsection{Experimental Setup}
% TODO: Dataset, metrics, baselines

\subsection{RQ1: Annotation Suggestion Effectiveness}
% TODO: Results

\subsection{RQ2: Program Repair Effectiveness}
% TODO: Results

\subsection{RQ3: LLM Comparison}
% TODO: Results (if applicable)

\subsection{Threats to Validity}
% TODO: Internal, external, construct validity

%% ============================================================
\section{Discussion}\label{sec:discussion}
%% ============================================================

% TODO: Limitations and future work

%% ============================================================
\section{Related Work}\label{sec:related}
%% ============================================================

\subsection{Intent Specification for Smart Contracts}
% TODO: IntentChecker, Scribble, other annotation approaches

\subsection{LLM-Based Program Repair}
% TODO: APR with LLMs, code generation

\subsection{LLM for Smart Contract Security}
% TODO: GPTScan, other LLM-based vulnerability detection

%% ============================================================
\section{Conclusion}\label{sec:conclusion}
%% ============================================================

% TODO: Summary and future work

%% ============================================================
%% Bibliography
%% ============================================================
\nocite{*}  % 모든 참고문헌을 인용 없이 포함
\bibliographystyle{ACM-Reference-Format}
\bibliography{references}

\end{document}
